\section{The Symmetric Case}
\label{sec:symmetric}

\subsection{Average protection and relative treatment}
\label{sec:decomposition}

Restrict attention to tariff changes imposed by $A$, $t=(t_j)_{j\neq A}$ with $t_j\equiv d\tau_{Aj}$, and let $G$ now denote the $(N-1)\times(N-1)$ restriction of the tariff-incidence matrix. Suppose $A$'s partners are \textbf{symmetric}: $PJP'=J$ and $PGP'=G$ for every permutation $P$ of the partners, as at free trade or under a common external tariff. Decompose
\begin{equation}
t=\bar t\,\bone+\tilde t,\qquad \bar t=\tfrac{1}{N-1}\bone't,\qquad \bone'\tilde t=0.
\label{eq:tsplit}
\end{equation}
We call $\bar t$ \textbf{average protection} and $\tilde t$ \textbf{relative treatment}. Any permutation-invariant matrix has the form $aI+b\bone\bone'$, so $\bone$ and every zero-sum vector are eigenvectors of both $G$ and $-J$:
\begin{equation*}
G\bone=g_S\bone,\quad Gv=g_Dv;\qquad (-J)\bone=j_S\bone,\quad (-J)v=j_Dv\qquad(\bone'v=0).
\end{equation*}

\begin{proposition}[Average protection and relative treatment]
\label{prop:decomposition}
Let partners be symmetric, $N\geq3$, and $j_S,j_D\neq0$. For any tariff vector $t$,
\begin{equation}
d\nu=\lambda_A\,\bar t\,\bone+\lambda_R\,\tilde t,\qquad \lambda_A=g_S/j_S,\quad \lambda_R=g_D/j_D.
\label{eq:decomp}
\end{equation}
Hence, with $d\nu^E_A\equiv\frac{1}{N-1}\sum_{j\neq A}d\nu_j$ (the NEER at symmetric weights),
\begin{equation}
d\nu^E_A=\lambda_A\bar t,\qquad d\nu_j-d\nu^E_A=\lambda_R\tilde t_j,\qquad d\log e_{jk}=\lambda_R(t_k-t_j).
\label{eq:decompparts}
\end{equation}
Moreover, MML is equivalent to $j_S>0$ and $j_D\geq j_S$.
\end{proposition}

\textbf{Average protection alone moves the NEER; relative treatment alone moves bilateral rates around it and the cross rates among partners.} The geometry is simple. The response vector lives in the plane spanned by the common direction $\bone$ and the zero-sum directions, and each component is scaled by its own coefficient (\cref{fig:geometry}).

\begin{figure}[t]
\centering
\begin{subfigure}[t]{0.45\textwidth}\centering
\includegraphics[width=\textwidth]{fig2a_geometry.pdf}
\caption{$\rho=2.2$: $\rho_D<\rho<3\rho_D$}
\label{fig:geometry-a}
\end{subfigure}\hfill
\begin{subfigure}[t]{0.45\textwidth}\centering
\includegraphics[width=\textwidth]{fig2b_geometry.pdf}
\caption{$\rho=5$: $\rho>3\rho_D$}
\label{fig:geometry-b}
\end{subfigure}
\caption{The geometry of average protection and relative treatment ($N=3$). A unilateral tariff $t=(1,0)=(\tfrac12,\tfrac12)+(\tfrac12,-\tfrac12)$, with $\alpha_D=0.7$, $\alpha_T=0.6$, $\eta=1.5$ ($\rho_D=1.17$, $3\rho_D=3.51$). Blue: common component $\lambda_A\bar t\,\bone$; solid orange: relative component $\lambda_R\tilde t$; black: equilibrium $d\nu$. The dashed orange arrow scales the relative component by $N$ to show $F/j_S$. In (a), $\rho_D<\rho<3\rho_D$: $C$'s trade balance improves at unchanged exchange rates ($F_C>0$), yet $A$ still appreciates against $C$. In (b), $\rho>3\rho_D$: $A$ depreciates against $C$. In both, the NEER appreciates. Computed from \cref{thm:main}.}
\label{fig:geometry}
\end{figure}

\subsection{A nested-CES economy}
\label{sec:model}

Each country $i\in\{1,\dots,N\}$ is endowed with $L_i$ units of labor, its size, which is mobile across sectors. It produces a country-specific tradable and a non-tradable with labor only, $Y_{Ti}=A_{Ti}L_{Ti}$ and $Y_{Ni}=A_{Ni}L_{Ni}$, with $L_{Ti}+L_{Ni}=L_i$. Normalizing each tradable's producer price to one in its own currency, $w_i=A_{Ti}$ and $P_{Ni}=A_{Ti}/A_{Ni}$. Relative international producer prices therefore move only through exchange rates, and \textbf{the bilateral exchange rate equals the bilateral terms of trade and the relative wage in a common currency.}

Preferences are
\begin{equation}
C_i=C_{Ti}^{\alpha_T}C_{Ni}^{1-\alpha_T},\quad
C_{Ti}=\Big[\alpha_D^{1/\eta}C_{Tii}^{\frac{\eta-1}{\eta}}+(1-\alpha_D)^{1/\eta}M_i^{\frac{\eta-1}{\eta}}\Big]^{\frac{\eta}{\eta-1}},\quad
M_i=\Big[\textstyle\sum_{j\neq i}\beta_{ij}^{1/\rho}C_{Tji}^{\frac{\rho-1}{\rho}}\Big]^{\frac{\rho}{\rho-1}}.
\label{eq:prefs}
\end{equation}
Here $\eta$ governs substitution between domestic tradables and the import composite, and $\rho$ substitution across foreign origins; the latter is inactive when $N=2$. Country $i$ levies $\tau_{ij}$ on imports from $j$, so the consumer price is $p_{ij}=(1+\tau_{ij})e_{ij}$, and the import and tradable price indices are $P_{Mi}=[\sum_{j\neq i}\beta_{ij}p_{ij}^{1-\rho}]^{1/(1-\rho)}$ and $P_{Ti}=[\alpha_D+(1-\alpha_D)P_{Mi}^{1-\eta}]^{1/(1-\eta)}$. Tariff revenue is rebated lump-sum, so income is
\begin{equation}
I_i=\frac{w_iL_i}{1-\sum_{j\neq i}\frac{\tau_{ij}}{1+\tau_{ij}}s_{ij}},
\label{eq:income}
\end{equation}
with $s_{ij}$ the realized tariff-inclusive expenditure share of income on $j$'s good. Trade balances are measured at border prices,
\begin{equation}
\TB_i=\sum_{k\neq i}f_{ki}-\sum_{j\neq i}f_{ij},\qquad f_{ij}=\frac{s_{ij}\tilde I_i}{1+\tau_{ij}},
\label{eq:tb}
\end{equation}
with $\tilde I_i=e_{Ai}I_i$ income in $A$'s currency. The NEER uses baseline bilateral trade weights, $E_i=\prod_{j\neq i}e_{ij}^{w_{ij}}$ with $w_{ij}=(f_{ij}+f_{ji})/\sum_k(f_{ik}+f_{ki})$; throughout, ``effective exchange rate'' refers to this trade-weighted NEER. The real exchange rate is $q_{ij}=e_{ij}P_j/P_i$, with $P_i$ the consumer price index.

The \textbf{symmetric benchmark} is free trade with $\beta_{ij}=1/(N-1)$, common parameters, and equal endowments and productivities. It retains home bias $\alpha_D$. Write $m\equiv\alpha_T(1-\alpha_D)$ for the import share of income.

\paragraph{The two-country case.} With one foreign supplier, linearizing gives $d\TB_A=mM_2\,d\nu_B+m\rho_D\,d\tau$, so
\begin{equation}
\begin{gathered}
\frac{d\log e_{AB}}{d\tau}\bigg|_{\tau=0}=-\frac{\rho_D}{M_2},\\
\rho_D\equiv\underbrace{\alpha_D\eta}_{\text{domestic substitution}}+\underbrace{(1-\alpha_D)(1-\alpha_T)}_{\text{rebate / non-tradables}},\qquad
M_2\equiv1+2\alpha_D(\eta-1)=\varepsilon_X+\varepsilon_M-1.
\end{gathered}
\label{eq:twocountry}
\end{equation}
$\rho_D$ is the \textbf{domestic diversion margin}. It measures how strongly a tariff pulls expenditure toward the domestic economy, through substitution toward the domestic tradable and through the rebate's income effect on non-tradables. An exchange-rate movement generates no revenue, so it has no rebate term. This is why $\rho_D\neq M_2$.

\subsection{The main theorem}
\label{sec:theorem}

\begin{theorem}[Exchange-rate responses in the symmetric nested-CES model]
\label{thm:main}
At the symmetric benchmark define
\begin{equation}
M_N\equiv\frac{1}{N-1}\Big[\,1+\underbrace{N\alpha_D(\eta-1)}_{\text{domestic--import adjustment}}+\underbrace{(N-2)\rho}_{\text{foreign-origin adjustment}}\Big].
\label{eq:MN}
\end{equation}
MML holds if and only if $M_N>0$, and $\eta\geq1$ suffices. When it holds,
\begin{equation}
\lambda_A=-\frac{\rho_D}{M_N},\qquad \lambda_R=-\frac{\rho}{NM_N},\qquad d\nu_j=\lambda_A\bar t+\lambda_R\tilde t_j.
\label{eq:lambdas}
\end{equation}
\end{theorem}

$M_N$ measures how sensitive $A$'s trade balance with the rest of the world is to its effective exchange rate: $d\TB_A=mM_N\,d\nu^E_A$. When $N>2$ it includes $\rho$, because foreign buyers can substitute between $A$'s goods and other foreign goods. For $N=2$ it reduces to $M_2$. Average protection works as in the two-country case. Relative treatment is zero-sum, so to first order it leaves revenue, income and the tradable/non-tradable split unchanged, and its disturbance is governed by $\rho$ alone.

The factor $1/N$ in $\lambda_R$ is central. A common exchange-rate move shifts only partners' trade with $A$. A relative move also shifts partners' trade with \emph{each other}. At the symmetric benchmark $A$ accounts for a share $1/N$ of the trade on which relative movements act, so $j_D=Nj_S$. \textbf{Relative exchange-rate adjustment is $N$ times as effective as common adjustment}, and a given relative disturbance requires only $1/N$ as much exchange-rate movement.

\subsection{A unilateral tariff}
\label{sec:unilateral}

Let $A$ raise its tariff on $B$ alone, $t=d\tau\,e_B$. Then $\bar t=d\tau/(N-1)$, $\tilde t_B=\frac{N-2}{N-1}d\tau$ and $\tilde t_C=-\frac{d\tau}{N-1}$ for every bystander $C$.

\begin{proposition}[Unilateral tariff]
\label{prop:unilateral}
At the symmetric benchmark under MML,
\begin{equation}
\frac{d\nu^E_A}{d\tau}=-\frac{\rho_D}{(N-1)M_N}<0,\qquad
\frac{d\nu_B}{d\tau}=-\frac{N\rho_D+(N-2)\rho}{N(N-1)M_N}<0,\qquad
\frac{d\nu_C}{d\tau}=\frac{\rho-N\rho_D}{N(N-1)M_N}.
\label{eq:unilateral}
\end{equation}
\begin{enumerate}[label=(\roman*)]
\item $A$'s NEER appreciates for every $N$ and $\rho$.
\item $A$ appreciates against the target.
\item $A$ \textbf{depreciates against each untariffed bystander if and only if $\rho>N\rho_D$.}
\item The same threshold applies when $A$ tariffs any $k\leq N-2$ partners. More generally, $A$ depreciates against a partner $j$ facing a below-average tariff increase if and only if $t_j/\bar t<1-N\rho_D/\rho$, which requires $\rho>N\rho_D$.
\end{enumerate}
\end{proposition}

Against the target, average protection and relative treatment reinforce each other. Against a bystander they pull apart. Average protection shifts expenditure toward the domestic economy (appreciation), while the bystander's favorable relative treatment redirects expenditure toward it (depreciation).

\begin{corollary}[Why a single elasticity cannot deliver reversal]
\label{cor:single}
With home bias $\alpha_D\geq1/N$,
\begin{equation}
N\rho_D-\eta=\eta(N\alpha_D-1)+N(1-\alpha_D)(1-\alpha_T)\geq0,
\label{eq:single}
\end{equation}
so $\rho>\eta$ is \textbf{necessary} for a bystander depreciation. In any specification with a single elasticity across all tradable varieties ($\rho=\eta=\sigma$), $A$ appreciates against every bystander at every $\sigma$. At $\alpha_D=1/N$ the response is $-(1-\alpha_T)/(N\sigma)$, which approaches zero from below but never crosses it.
\end{corollary}

This explains why the natural first extension of the two-country model, a third country under a single CES elasticity, has not surfaced the reversal. The reversal requires that cross-origin substitution exceed home--import substitution, an ordering supported by the macro--micro elasticity evidence \citep{feenstra2018,li2022}. At US-like parameters ($\alpha_D=0.8$, $\alpha_T=0.4$, $\eta=1.5$, so $\rho_D=1.32$), $N\rho_D=3.96$ for $N=3$.

\paragraph{Disturbance versus equilibrium.} At unchanged exchange rates the bystander's trade balance improves if and only if $\rho>\rho_D$:
\begin{equation}
F_C=\frac{m(\rho-\rho_D)}{(N-1)^2}.
\label{eq:FC}
\end{equation}
The equilibrium bystander depreciation requires the stronger condition $\rho>N\rho_D$. The gap is the $1/N$ effectiveness of relative adjustment. For $\rho_D<\rho<N\rho_D$, $C$ gains at the old exchange rates yet ends up with a weaker currency against $A$ (\cref{fig:geometry-a}). \Cref{fig:unilateral} plots the three responses against $\rho$.

\begin{figure}[t]
\centering
\begin{subfigure}[t]{0.45\textwidth}\centering
\includegraphics[width=\textwidth]{fig3a_unilateral.pdf}
\caption{$N=3$}
\end{subfigure}\hfill
\begin{subfigure}[t]{0.45\textwidth}\centering
\includegraphics[width=\textwidth]{fig3b_unilateral.pdf}
\caption{$N=5$}
\end{subfigure}
\caption{Bilateral and effective responses to a unilateral tariff, per unit $d\tau$, US-like parameters ($\alpha_D=0.8$, $\alpha_T=0.4$, $\eta=1.5$). $d\nu_B$ and $d\nu^E_A$ are negative for every $\rho$; $d\nu_C$ crosses zero at $N\rho_D$ ($3.96$ for $N=3$, $6.60$ for $N=5$), and the shaded region is $\rho>N\rho_D$. The grey dotted line is the single-elasticity case $\rho=\eta=\sigma$, plotting $d\nu_C$ against $\sigma$; it never crosses zero (\cref{cor:single}). Closed forms of \cref{prop:unilateral}, verified against the linear solver.}
\label{fig:unilateral}
\end{figure}

\subsection{Retaliation and trade wars}
\label{sec:retaliation}

Since the analysis is local, simultaneous tariff changes by several countries add up.

\begin{proposition}[Retaliation and trade wars]
\label{prop:retaliation}
At the symmetric benchmark with $N\geq3$ under MML:
\begin{enumerate}[label=(\roman*)]
\item \textbf{Bilateral retaliation} ($A$ tariffs $B$ and $B$ tariffs $A$ by $d\tau$): $d\nu_B=0$, and
\begin{equation}
\frac{d\nu_C}{d\tau}=\frac{\rho-\rho_D}{(N-1)M_N},\qquad
\frac{d\nu^E_A}{d\tau}=\frac{N-2}{N-1}\cdot\frac{\rho-\rho_D}{(N-1)M_N}.
\label{eq:retaliation}
\end{equation}
Both belligerents depreciate against every bystander, and in effective terms, \textbf{if and only if $\rho>\rho_D$.} Cross rates among bystanders are unchanged.
\item \textbf{Multilateral trade war} ($A$ tariffs every partner, and every partner retaliates against $A$): $d\nu_j=d\nu^E_A=\frac{(N-2)(\rho-\rho_D)}{(N-1)M_N}\,d\tau$ for all $j$, so $A$ depreciates against everyone if and only if $\rho>\rho_D$.
\end{enumerate}
\end{proposition}

$B$'s retaliation adds a second relative-treatment effect in the bystander's favor, which offsets the $1/N$ dilution. The threshold drops from $N\rho_D$ to $\rho_D$, the point at which the bystander's trade balance improves at unchanged exchange rates. At US-like parameters $\rho_D=1.32$, below essentially all cross-origin estimates. \textbf{Retaliation therefore makes multilateral depreciation of the tariffing country's currency the generic outcome, while leaving the belligerents' bilateral rate unchanged.} This is the pattern documented by \citet{ostry2025}, including the near-stability of the dollar against the renminbi in tit-for-tat episodes. Here it arises with flexible prices, balanced trade and no financial channel.

The two experiments differ in where adjustment falls. Under bilateral retaliation the belligerents' rate is fixed and adjustment runs against bystanders. In a multilateral war there are no bystanders, and all of $A$'s rates move together (\cref{fig:configurations}).

\begin{figure}[t]
\centering
\includegraphics[width=0.9\textwidth]{fig4_configurations.pdf}
\caption{Configuration dependence, $N=3$, US-like parameters. NEER responses $d\nu^E_A/d\tau$ under a unilateral tariff (always negative), bilateral retaliation and a multilateral war (both crossing zero at $\rho_D=1.32$), and the bilateral response against a bystander under a unilateral tariff (crossing at $3\rho_D=3.96$) and under retaliation (crossing at $\rho_D$). At $N=3$ the multilateral-war response coincides with the bystander response under bilateral retaliation, so those two lines overlap. Shaded: the central empirical range for cross-origin elasticities, $\rho\in[2,4]$.}
\label{fig:configurations}
\end{figure}

\subsection{The three-country representation}
\label{sec:composite}

The $N$-country results can be read through a three-country economy, which we use as the workhorse for intuition. When the $N-2$ bystanders are symmetric and the tariff shock preserves that symmetry, no relative-price movement within the bystander bloc is activated. The bloc then has a single active direction, and the tariff experiment reduces exactly to three regions: $A$, $B$ and a composite bystander $C$. We now show that this composite is itself a nested-CES country, and we identify its parameters.

\begin{proposition}[Exact three-country reduction]
\label{prop:composite}
At the symmetric benchmark, consider any tariff change among $A$, $B$ and the bystanders that is common across bystanders: tariffs by $A$ or $B$ on all bystanders, by all bystanders on $A$ or $B$, or between $A$ and $B$. To first order, the responses of $\nu_B$, of every bystander rate $\nu_C$ and of $A$'s NEER coincide with those of a three-country nested-CES economy $\{A,B,C\}$ in which:
\begin{itemize}
\item $C$ has income $y_C=N-2$;
\item $A$ and $B$ place import weight $\frac{N-2}{N-1}$ on $C$ and $\frac{1}{N-1}$ on each other;
\item $C$ places import weight $\tfrac12$ on each of $A$ and $B$, with cross-origin elasticity $\rho$;
\item $C$'s home share of tradables is $h_C=\alpha_D+(1-\alpha_D)\frac{N-3}{N-1}$, equivalently $(N-1)h_C=N-3+2\alpha_D$;
\item $C$'s home--import elasticity is
\begin{equation}
\boxed{\;\eta_C=\frac{(N-3)\rho+2\alpha_D\,\eta}{N-3+2\alpha_D}.\;}
\label{eq:etaC}
\end{equation}
\end{itemize}
\end{proposition}

\begin{proof}[Proof sketch]
All bystander goods have the same price change, so CES aggregation of $A$'s and $B$'s demand over bystander goods is exact. Trade within the bloc cancels from the bloc's trade balance; what remains is the bloc's demand for $A$'s and $B$'s goods. Measure prices in bystander currency, so bloc goods have $dp=0$, and write $S=dp_A+dp_B$ and $D=dp_A-dp_B$. A bystander's import price index moves by $S/(N-1)$, since $A$ and $B$ are 2 of its $N-1$ import origins. Hence
\begin{equation*}
d\log s_{cA}=(1-\rho)\tfrac{D}{2}+\Big[(1-\rho)\tfrac{N-3}{2(N-1)}+(1-\eta)\tfrac{\alpha_D}{N-1}\Big]S.
\end{equation*}
In the composite the import index moves by $S/2$, so $d\log s_{CA}=(1-\rho)\frac D2+(1-\eta_C)h_C\frac S2$. The $D$ terms coincide. Matching the $S$ terms gives $(1-\eta_C)(N-1)h_C=(1-\rho)(N-3)+2\alpha_D(1-\eta)$, which is \cref{eq:etaC}. The full proof is in \cref{app:proofs}.
\end{proof}

The reduction is first-order. At finite tariffs a CES with constant $\eta_C$ does not replicate the bloc.

\paragraph{Interpretation.} The composite bystander differs from an ordinary large country in one respect. Most of each bystander's trade is with \emph{other bystanders}, whose goods sit in the same $\rho$-nest as $A$'s and $B$'s. Once collapsed, this internal trade becomes the composite's ``home'' trade, and it is substitutable at a blend of $\rho$ and $\eta$. When $\rho>\eta$, $\eta_C$ rises toward $\rho$ as $N$ grows. The bystander bloc becomes inward-looking ($h_C\to1$) and highly substitutable.

\paragraph{The reduced system.} In the composite economy the $2\times2$ system for $(\nu_B,\nu_C)$ is, with $c=mM_N/(N-1)$,
\begin{equation}
-J=c\begin{pmatrix}N-1&-(N-2)\\-(N-2)&2(N-2)\end{pmatrix},\qquad \det(-J)=c^2N(N-2),
\label{eq:reducedJ}
\end{equation}
and
\begin{equation}
d\nu_C=\frac{F_B+(N-1)F_c}{cN},\qquad d\nu_B=\frac{2F_B+(N-2)F_c}{cN},
\label{eq:reducedsol}
\end{equation}
where $F_B$ is the target's disturbance and $F_c$ each bystander's ($F_C=(N-2)F_c$). Under a unilateral tariff, $F_B=-m[\rho_D+(N-2)\rho]/(N-1)^2<0$ and $F_c=m(\rho-\rho_D)/(N-1)^2$. \Cref{eq:reducedsol} recovers \cref{eq:unilateral} and gives the bystander response a two-force reading:
\begin{itemize}
\item \textbf{diversion}, $(N-1)F_c$: expenditure diverted toward bystanders improves their position and strengthens their currencies;
\item \textbf{redirection}, $F_B$: the target loses access to $A$'s market and its currency must fall. This makes the target's goods cheaper for bystanders and drags bystander currencies down with it.
\end{itemize}
$A$ depreciates against the bystander if and only if diversion outweighs redirection, $(N-1)(\rho-\rho_D)>\rho_D+(N-2)\rho$, which is $\rho>N\rho_D$. These are the ``trade diversion'' and ``trade redirection'' forces of earlier versions of this paper, now with exact weights.

\paragraph{Why $N$ raises the threshold, and why size alone does not.} One might expect that adding countries is equivalent to enlarging the bystander. It is not. Hold the composite's flows fixed at their $N$-country values but give it an ordinary home--import elasticity $\eta$ instead of $\eta_C$. The bystander threshold then \emph{falls} slightly with $N$, consistent with the general result (\cref{app:asym-primitives}) that a larger bystander lowers the threshold. The true threshold $N\rho_D$ rises steeply (\cref{tab:composite,fig:composite}). What drives the $N$-dependence is the bloc's internal, highly substitutable trade, which makes relative exchange-rate adjustment $N$ times as effective as common adjustment and dilutes $A$'s footprint to $1/N$.

\begin{table}[t]
\centering
\caption{$N$ versus bystander size}
\label{tab:composite}
\begin{tabular}{ccccc}
\toprule
 & & \multicolumn{2}{c}{Composite bystander threshold} & \\
\cmidrule(lr){3-4}
$N$ & $N\rho_D$ (true) & with $\eta_C$ \cref{eq:etaC} & with ordinary $\eta$ & $\eta_C$ at $\rho=4$ \\
\midrule
\tabinput{tables/tab1_composite.tex}
\bottomrule
\end{tabular}

\vspace{4pt}
\begin{minipage}{0.92\textwidth}\footnotesize
\textit{Notes:} US-like parameters ($\alpha_D=0.8$, $\alpha_T=0.4$, $\eta=1.5$). Thresholds are the values of $\rho$ at which the bystander response to a unilateral tariff changes sign in the composite economy of \cref{prop:composite}, computed with the linear system of \cref{app:linear}. The ``ordinary $\eta$'' column keeps the composite's flows but replaces $\eta_C$ by $\eta$.
\end{minipage}
\end{table}

\begin{figure}[t]
\centering
\begin{subfigure}[t]{0.45\textwidth}\centering
\includegraphics[width=\textwidth]{fig5a_composite_threshold.pdf}
\caption{Bystander threshold}
\end{subfigure}\hfill
\begin{subfigure}[t]{0.45\textwidth}\centering
\includegraphics[width=\textwidth]{fig5b_eta_C.pdf}
\caption{Effective elasticity $\eta_C$}
\end{subfigure}
\caption{The composite bystander, US-like parameters. (a) The true threshold $N\rho_D$ for $N=3,\dots,10$, which the composite with $\eta_C$ reproduces exactly, and the threshold of the ``ordinary-$\eta$'' composite, which is nearly flat and slightly declining. (b) $\eta_C$ from \cref{eq:etaC} against $N$ for $\rho\in\{2,4,8\}$, converging toward $\rho$ (thin horizontal lines); the grey line is $\eta=1.5$.}
\label{fig:composite}
\end{figure}

\subsection{Real exchange rates}
\label{sec:real}

Real rates add a price-index wedge: the tariff raises $A$'s consumer prices directly. With $\bar t$ the average tariff and partners imposing no tariffs,
\begin{equation}
d\log q_{Aj}=\Big(1-\tfrac{mN}{N-1}\Big)d\nu_j-m\bar t,\qquad
dq^E_A=\Big[\Big(1-\tfrac{mN}{N-1}\Big)\lambda_A-m\Big]\bar t.
\label{eq:real}
\end{equation}

\begin{proposition}[Real exchange rates]
\label{prop:real}
Consider a unilateral tariff at the symmetric benchmark.
\begin{enumerate}[label=(\roman*)]
\item If $m<1/N$ and MML holds at the relevant $\rho$, the real bilateral threshold exceeds the nominal one:
\begin{equation*}
\rho^*_q(N)=\frac{[(N-1)-mN]N\rho_D+mN[N\alpha_D(\eta-1)+1]}{(N-1)(1-mN)}>N\rho_D.
\end{equation*}
\item If $m<(N-1)/N$, the REER appreciates.
\end{enumerate}
\end{proposition}

At US-like parameters and $N=3$, $\rho^*_q=4.93$, against $3\rho_D=3.96$.

\subsection{Relation to the competing-exporter model}
\label{sec:pta}

In our model class the bystander's exchange rate is its terms of trade: $\ToT_{CA}=e_{AC}$. The response in \cref{eq:unilateral} is therefore the third-country terms-of-trade effect of a discriminatory tariff, the central object of the preferential-trade-agreement literature since \citet{viner1950} and \citet{mundell1964}. In the canonical competing-exporter model, an importer $A$ that does not produce the good buys it from $B$ and $C$ under quasi-linear preferences. With $q_j$ exporter $j$'s price, $E_j$ its export supply and $M_A$ the importer's demand, a tariff on $B$ then \emph{always} improves $C$'s terms of trade, since
\begin{equation*}
\frac{dq_C}{dt_B}=\frac{E'_B}{E'_B+E'_C-M'_A}>0,
\end{equation*}
and preferential liberalization toward $B$ always hurts the excluded country. In our model the same effects hold \textbf{only if $\rho>N\rho_D$}. The difference is that $A$ produces a competing tradable. Part of the demand diverted from $B$ goes to $A$'s own good, along with the rebate's income effect, so the third-country effect survives only when cross-origin substitution is strong relative to domestic diversion. The two environments are not formally nested. The competing-exporter model can be read as switching off $\rho_D$ in order to isolate substitution across foreign suppliers. \Cref{app:asym-tariff} shows that at a positive common external tariff the threshold rises, and that a discrete preferential cut requires a higher $\rho$ than the local threshold suggests.
